% chapters/design/ocr_pipeline.tex
% Sección 4.5: Pipeline de ingesta y OCR (~1.5-2 páginas)
% ============================================================================
%
% GUÍA: Describir el flujo completo desde que un PDF llega por SFTP hasta
%       que se crea una instancia de examen vinculada a un alumno.
%       Referencia a \ref{rf:ocr} y \ref{rnf:maint:ocr}.
%
%   --- DIAGRAMA 7: Diagrama de flujo del pipeline OCR ---
%   Este es el diagrama de flujo que te recomendó el tutor. Mostrar:
%     1. SFTP watcher detecta nuevo PDF
%     2. Se encola tarea en Celery (Redis)
%     3. Worker recoge la tarea
%     4. Se almacena el PDF original en MinIO
%     5. Se extraen las páginas/regiones de interés
%     6. Se invoca el plugin OCR sobre las zonas configuradas
%     7. Se obtiene texto + confianza de cada zona
%     8. Se identifica al alumno por similitud (NIA, DNI, nombre)
%     9. Decisión: ¿identificación exitosa?
%        - Sí → se crea ExamInstance vinculada al alumno
%        - No → se crea incidencia para resolución manual
%    10. Se registra en auditoría
%
%   \begin{figure}{fig:ocr_flow}{Diagrama de flujo del \textit{pipeline} de ingesta y \ac{ocr}}
%     \image{}{}{ocr_flow_diagram}
%   \end{figure}
%
%   TODO: Crear la imagen figures/ocr_flow_diagram.png (o .pdf).
%         Herramientas: draw.io, Mermaid, Excalidraw.
%
%   --- DIAGRAMA 8: Diagrama de secuencia del procesamiento OCR ---
%   Mostrar la interacción entre componentes: SFTPWatcher → Redis →
%   CeleryWorker → MinIO → OCRPlugin → PostgreSQL.
%
%   \begin{figure}{fig:ocr_sequence}{Diagrama de secuencia del procesamiento \ac{ocr}}
%     \image{}{}{ocr_sequence_diagram}
%   \end{figure}
%
%   TODO: Crear la imagen figures/ocr_sequence_diagram.png (o .pdf).
%
%   --- CONTENIDO TEXTUAL ---
%   1. Describir cada etapa del pipeline.
%   2. Explicar la arquitectura pluggable del motor OCR:
%      - Interfaz: entrada (imagen) → salida (texto, confianza).
%      - Implementación concreta: Tesseract / EasyOCR / etc.
%      - TODO: Especificar qué motor OCR usas.
%   3. Algoritmo de identificación por similitud:
%      - Campos usados: NIA, DNI, nombre.
%      - Métrica de similitud (Levenshtein, fuzzy matching, etc.).
%      - Umbral de confianza.
%   4. Gestión de incidencias / errores del OCR.
%   5. Idempotencia y reintentos (si un worker falla).
% ============================================================================

% TODO: Redactar la sección del pipeline OCR y crear los diagramas.
